% Section 12.1, from lectures/L12/appendix/a_meta_prev.md.

\section{Designing \code{meta_prev}}
\label{c12:sec:metaprev}

\code{rx_bus} arrives from a transceiver on another chip, and \code{reset_n} from a button or a
supervisor, so nothing in this FPGA has ever controlled either. Sampling one of them straight into
\code{clock}'s domain can leave a flip-flop metastable, which is the problem \chapref{c4:ch} solved
with two flip-flops in series. You have built this module before, as \code{sync} in
Exercise~\ref{c5:ex:sync}; this is the same idea, written once, properly, to be kept in this design.

\Secref{c11:sec:sync} already settled \emph{that} the controller synchronises its own two
asynchronous inputs, and why: they are asynchronous whoever instantiates \code{can_controller}, so
the controller owns them and an integrator connects both straight to pins without thinking about it
further. This section builds the module that does it. It is the only module in the whole book that
contains nothing at all about CAN, which is exactly why it is reusable enough to end up with three
instances by \chapref{c19:ch}.

\subsection{Interface}
\label{c12:sub:mpiface}

\modulefig[0.62]{lectures/L12/appendix/images/meta_prev.png}

\begin{booktable}{@{}p{5em}p{4.4em}p{11.4em}L@{}}
  \thead{Port / generic} & \thead{Dir.} & \thead{Type} & \thead{Meaning} \\
  \midrule
  \code{WIDTH} & generic & \code{natural := 1} & Number of independent bits to synchronise. \\
  \code{clock} & in & \code{std_logic} & 50 MHz system clock, the domain being crossed \emph{into}. \\
  \code{async_in} & in & \code{std_logic_vector(WIDTH-1 downto 0)} & The asynchronous input or
    inputs. \\
  \code{sync_out} & out & \code{std_logic_vector(WIDTH-1 downto 0)} & The same bits, safe to use in
    this domain. \\
\end{booktable}

\code{meta_prev_tb} binds its ports positionally, as everything in this book does, so the order and
the types above have to match exactly. \textbf{The generic map is positional too}, so \code{WIDTH} has
to be the module's only generic (or at least its first). The testbench instantiates the module twice,
once with the default width and once with width 4, so that both the default and the generic get
exercised. It is also the only module here that reads no package at all, so it needs neither
\code{can_def} nor anything else.

\subsection{Behaviour}
\label{c12:sub:mpbehav}

Two flip-flops per bit, in series, on \code{clock}'s rising edge: the first can go metastable, and
the second gives it a whole clock period to settle before anything downstream sees it.
\code{sync_out} therefore follows \code{async_in} after exactly two rising edges.

Three things differ from the \code{sync} module you wrote in Exercise~\ref{c5:ex:sync}, and all three
are deliberate:
\begin{itemize}
  \item \textbf{No reset port.} \code{sync} took a \code{reset_s2_n}. This one takes none: after two
    clock edges a synchroniser holds real data regardless of what it started from. Dropping the port
    is also what lets the \emph{same} module synchronise a reset, without the circularity of a reset
    synchroniser needing a synchronised reset. The testbench clocks the initial level through before
    checking anything, for exactly that reason.
  \item \textbf{No \code{PRESET} generic.} \Chapref{c5:ch} argued at length that the safe reset value
    depends on what the signal \emph{means}, and that only the instantiating design knows that. That
    argument was about the two cycles after reset, and it does not survive dropping the reset:
    without a reset branch there is no value to preset. What replaces it is the caller's obligation
    to ignore \code{sync_out} during the first two edges, which is the same obligation expressed
    differently.
  \item \textbf{\code{SIZE} is called \code{WIDTH}.} The same generic, the same job, renamed to the
    name every subblock in this project uses.
\end{itemize}

This design instantiates it three times, all with \code{WIDTH = 1}: twice inside
\code{can_controller}, for \code{reset_n} and for \code{rx_bus}, and once at \code{can_spi_node}'s
top level for its own reset (\chapref{c19:ch}). One module, three instances, no copies. The generic
still earns its place \textemdash{} the testbench exercises \code{WIDTH = 4}, and a synchroniser
written for one bit is a synchroniser you copy the next time you need four.

\begin{warning}
A caveat to mention now, since \chapref{c19:ch} leans on it: a synchroniser handles metastability,
not \emph{skew}. The second flip-flop makes it overwhelmingly unlikely that a metastable value
propagates (never impossible; the probability falls exponentially with settling time, which is why
two flip-flops are the accepted trade-off), but every bit still crosses independently, so two bits
changing together can land one clock apart. That is why a synchroniser is not a general answer for a
bus whose bits have to stay together \textemdash{} and it is why \code{can_spi_node} synchronises its
reset with this module but leaves the SPI lines to \code{spi_slave}, which frames them, rather than
widening a \code{meta_prev} across them.
\end{warning}

\subsection{What the testbench pins down}
\label{c12:sub:mptb}

\begin{itemize}
  \item \code{sync_out} follows \code{async_in} after exactly two rising edges: not one (a single
    flip-flop is no synchroniser) and not three.
  \item The same holds from high to low as from low to high.
  \item With \code{WIDTH = 4} all four bits arrive together, two edges after the input changed.
  \item A stable input leaves \code{sync_out} stable.
\end{itemize}

This is the first testbench in the project that actually \emph{runs} rather than merely analyses, so
it is also where the third GHDL step arrives. \Chapref{c11:ch} analysed and elaborated; see
appendix~\ref{app:sim} for what \code{ghdl -r} adds and why \code{--assert-level=error} is not
optional.

\subsection{Putting it into \code{can_controller}}
\label{c12:sub:mpwire}

\code{meta_prev}'s ports are vectors even at width 1, so a scalar port needs a one-element signal on
each side. In the architecture's declarative part (between \code{is} and \code{begin}):

\begin{vhdlcode}
    -- meta_prev, synchronizing the two inputs that arrive from off-chip.
    signal reset_n_v, reset_s2_n_v : std_logic_vector(0 downto 0);
    signal rx_bus_v,  rx_bus_s2_v  : std_logic_vector(0 downto 0);
    signal reset_s2_n, rx_bus_s2   : std_logic;
\end{vhdlcode}

and after \code{begin}:

\begin{vhdlcode}
    reset_n_v(0) <= reset_n;
    rx_bus_v(0)  <= rx_bus;

    reset_s2_n <= reset_s2_n_v(0);
    rx_bus_s2  <= rx_bus_s2_v(0);

    reset_sync: entity work.meta_prev
        port map(clock, reset_n_v, reset_s2_n_v);

    rx_bus_sync: entity work.meta_prev
        port map(clock, rx_bus_v, rx_bus_s2_v);
\end{vhdlcode}

Two instances rather than one with \code{WIDTH => 2}. Both bits would ride along together perfectly
well, since they are independent and skew between them means nothing, but one named instance per
concern reads better on the diagram, and later making the reset assert asynchronously would then tear
nothing up.

From now on, \textbf{every block takes \code{reset_s2_n}, never \code{reset_n}, and every block that
reads the bus reads \code{rx_bus_s2}, never \code{rx_bus}}. That is the whole reason these two
instances exist, it applies to \code{bit_timer} in the very next section, and it is easy to forget
again when \code{rx_shift_reg} is wired in in \chapref{c15:ch}.
